\section{Исходные данные}
\label{sec:data}

Для экспериментов использованы три синтетических RGB-изображения размером $256 \times 256$ пикселей, различающиеся пространственной структурой:
\begin{itemize}
	\item \texttt{smooth.bmp}~--- плавные градиенты и низкочастотные переходы цвета;
	\item \texttt{contours.bmp}~--- крупные однородные области с резкими границами;
	\item \texttt{texture.bmp}~--- мелкая высокочастотная текстура и шум.
\end{itemize}
Изображения созданы скриптом \texttt{generate\_samples.py} с фиксированным seed; при необходимости они могут быть заменены любыми BMP того же или большего размера.

Параметры разбиения и характеристики выборки приведены в таблице~\ref{tab:data}. В основных сериях блоки не перекрываются; для демонстрации влияния перекрытия используется шаг $s_h=s_w=4$.

\begin{table}[H]
	\centering
	\caption{Исходные изображения и параметры выборки}
	\label{tab:data}
	\begin{tabular}{@{}lcccccc@{}}
		\hline
		Изображение & $h \times w$ & Каналов & $r \times m$ & $s_h \times s_w$ & $Q$ & $n$ \\
		\hline
		\texttt{smooth.bmp}   & $256\times256$ & 3 & $8\times8$ & $8\times8$ & 1024 & 192 \\
		\texttt{contours.bmp} & $256\times256$ & 3 & $8\times8$ & $8\times8$ & 1024 & 192 \\
		\texttt{texture.bmp}  & $256\times256$ & 3 & $8\times8$ & $8\times8$ & 1024 & 192 \\
		\hline
	\end{tabular}
\end{table}

Базовые параметры: $p=16$, $\beta=1$, $T_{\max}=80$, seed $=42$. При таком $p$ коэффициент по методике $Z \approx 10{,}1$, байтовый коэффициент $K \approx 1{,}26$, то есть оба показателя превышают единицу. В серии с перекрытием $s_h=s_w=4$ получается $Q=3969$, $Z\approx11{,}4$, $K\approx1{,}43$.

Среда выполнения: Python 3.14, NumPy, Pillow, Matplotlib; расчёты выполнены в одном потоке, время измерялось функцией \texttt{time.perf\_counter()}.
